\section{Review Against Related Work}
\label{sec:review}

\subsection{Where AstroFixxer differs}

Against AstroHopper and SkEye, AstroFixxer adds a way to measure the pointing instead of trusting the sensors after one alignment. Against StarSense Explorer and PiFinder, it needs no dock, mirror or separate computer; the phone already on the tube does the work. Unlike astrometry.net, the solver does not attempt blind solving at every scale. It uses the scale implied by the observer's equipment and the rough pointing from the sensors, which keeps the star file to 2.94 MB and a hinted solve to about half a second on a desktop CPU. The catalogue is larger than AstroHopper's and is fully offline, which matters at dark sites without mobile coverage. The Indian sky culture and the Hindi interface have no counterpart in the other free tools reviewed here.

\subsection{Limitations and what is not yet tested}

The evidence in Section~\ref{sec:system} comes from unit tests, synthetic images and a desktop test harness. None of it replaces a night at a telescope, and the following points remain open.
\begin{itemize}
\item The plate solver has only seen rendered star fields. Real phone photographs have non-Gaussian star images, compression artefacts, lens distortion beyond the 2\% tested, and noise that depends on the phone's processing. The detection thresholds were tuned on synthetic data.
\item Solving through an eyepiece depends on the phone recording stars of magnitude 9 to 10 in a short exposure. This has not been measured.
\item The revised alignment flow has only run in the desktop test harness, and the camera capture code is still under development. Neither has run on a phone.
\item The coordinate model omits nutation and aberration. At 28 arcseconds this is well inside an eyepiece field, but it would matter for setting-circle readouts on a precisely polar-aligned mount.
\item Sensor drift between alignments has not been measured on real phones or near real tubes.
\end{itemize}
A field protocol is already written: three phones (with and without a magnetometer), two telescopes, and ten Messier objects reached from alignment stars 5 to 20 degrees away, recording the time to target and the misses.
